\documentclass[11pt]{article}
\usepackage[margin=1in]{geometry}
\usepackage{amsmath,amssymb,amsthm}
\usepackage{mathtools}
\usepackage{enumitem}
\usepackage{hyperref}
\usepackage{cleveref}
\usepackage{booktabs}
\usepackage{tabularx}
\usepackage{microtype}
\usepackage{cite}
\usepackage{xurl}
\usepackage{path}
\hypersetup{colorlinks=true,linkcolor=blue,citecolor=blue,urlcolor=blue}
\newtheorem{definition}{Definition}
\newtheorem{lemma}{Lemma}
\title{Requestable Evidence Classes for Source-Only Releases\\\large Stable Names for Omitted-Support Families behind Packet Cuts and Public Request Contracts}
\author{Anonymous}
\date{Unified archive note v0.13 (2026-03-22; archive v615)}
\begin{document}
\maketitle

\begin{abstract}
A terse source-only paper can now say exactly which maintained packet/menu cut stands behind its ``supporting artifacts are available on request'' sentence. One repeat ambiguity remains: what \emph{kind} of omitted support is that cut promising? A referee may want bundle identity and validator companions, another may want the public-status lineage and closure slice, and another may want the compact-diff compare and replay slice. This note gives that vocabulary one home. It defines compact \emph{requestable evidence classes}: stable names for recurring omitted-support families that packet cuts and public request contracts may cite without re-enumerating the same paths in prose.
\end{abstract}

\paragraph{Non-redundancy note.}
Synthesis~6 owns the archive-wide source-only artifact policy, Synthesis~52 owns exact answer disclosure packets, Synthesis~54 owns request-class review menus, Synthesis~56 owns exact paper-facing public request contracts, Synthesis~58 owns exact completion criteria for when a promised cut has been fully satisfied, and Synthesis~17 owns the concrete worked instantiation.\cite{series:artifactpolicy,series:challengeanswerdisclosurepackets,series:challengeanswerreviewmenus,series:publicrequestcontracts,series:requestfulfillmentcertificates,series:workedexample}
This note owns only the reusable evidence-class vocabulary that those packet and contract notes may import when they need to say what omitted-support family a maintained cut is actually handing over. The public-anchor-family keys those classes mention remain imported contract vocabulary from Synthesis~56 rather than new family-basis ownership inside the evidence-class ledger. Exact completion criteria for when that handoff has been fully satisfied belong in Synthesis~58 rather than here, the within-release paper-level public status token for the overall promise belongs in Synthesis~59 rather than here, successor-release carryforward of that public token belongs in Synthesis~60 rather than here, the exact outward-facing notice sentence readers should see belongs in Synthesis~61 rather than here, the canonical sentence family plus slot realization rule behind that wrapper belong in Synthesis~62 rather than here, the request-side terminal stop belongs in Synthesis~73 rather than here, the paired answer/request citation discipline belongs in Synthesis~74 rather than here, and the stop-versus-widen/reopen rule belongs in Synthesis~75 rather than here.

\paragraph{Scope.}
This is not a new packet, menu, validator, or release stage.
It is a thin typing note saying that a maintained public-on-request handoff may carry one or more named evidence classes, each with a canonical path family and one public-anchor-family coverage claim imported from the companion public request contract.

\paragraph{Exports.}
This note exports (i) requestable evidence classes, (ii) a family-coverage rule, (iii) a packet-typing rule, (iv) a contract-coverage rule, (v) a shared-class reuse rule, and (vi) a shared-evidence-vocabulary rule saying that later papers may cite one stable evidence-class name when they need an auditable answer to ``what kind of omitted support is this source-only sentence promising?''

\paragraph{Later-terse-paper citation posture.}
Later terse papers should cite Synthesis~57 only when the live issue is the reusable omitted-support family typing itself: which stable evidence-class name characterizes one maintained handoff, whether that class is shared across more than one public-anchor family only as vocabulary, or which canonical support-path meaning later contracts, certificates, packets, routes, or spines are merely importing. Stop at Synthesis~56 when the live issue is the paper-facing promise that imports those class names. Widen to Synthesis~58 when the live issue is whether one family-scoped support cut has actually been fulfilled, widen to Synthesis~59--63 when the class basis is already fixed and the remaining issue is the paper-level status token, visible sentence, canonical notice family, or token-keyed family selector shown to readers, and stop at Synthesis~17 when the concrete maintained worked route is the live issue instead. Under the paired terminal citation discipline and paired cutover rule, later terse papers should usually stop at Synthesis~73 for the maintained request-side terminal stop, drop to Synthesis~17 when one concrete checked-answer/request-tail seam is wanted, and widen to Synthesis~74--75 only for the abstract paired-tail discipline or abstract stop-versus-widen judgment rather than reopen this vocabulary note.\cite{series:publicrequestcontracts,series:requestfulfillmentcertificates,series:publicrequeststatusenvelopes,series:publicrequestcarryforward,series:publicrequestnotices,series:publicrequestnoticenormalforms,series:publicrequestnoticeselections,series:requestsideterminalcitationnormalform,series:pairedterminalcitationdiscipline,series:pairedterminalcutoverrules,series:workedexample}

\paragraph{Within-release pre-wrapper route.}
This note is the second stop in the same four-stop within-release pre-wrapper route: one already-fixed paper-facing promise from Synthesis~56 narrows which evidence-class names can matter here, Synthesis~57 then says which omitted-support family each maintained handoff is typed as, Synthesis~58 decides whether each such family is complete, and Synthesis~59 compresses those family verdicts to one paper-level token before any visible notice is chosen. So this note owns family typing, not promise formation, not family completion, and not paper-level compression.\cite{series:publicrequestcontracts,series:requestfulfillmentcertificates,series:publicrequeststatusenvelopes}

\paragraph{One tiny route-scope drill.}
If two later terse papers agree on the same worked contract but disagree about whether the promised omitted support should be called a compact-diff family or a public-status-lineage family, the live issue belongs here. If they instead agree on the class name and disagree about whether that class has been fully discharged, they should widen immediately to Synthesis~58; if they agree on all family verdicts and only need the paper-level token, they should widen to Synthesis~59 instead of stretching this vocabulary note into a completion or status owner.\cite{series:workedexample}

\section{Why requestable evidence classes deserve their own note}
Archive-wide policy is intentionally broad.
Disclosure packets are intentionally answer-local.
Review menus are intentionally request-class-local.
Public request contracts are intentionally paper-local.
A short paper still benefits from one extra stable layer: not another packet, but a compact name for the support family that packet or contract is handing over.
Without that layer, one note says ``ask for the validator bundle,'' another says ``ask for the status lineage slice,'' and a third repeats path lists inline.
That is the same drift problem at a smaller granularity.

\begin{definition}[Requestable evidence class]
Fix one source-only public release and one public-anchor family $f$.
A \emph{requestable evidence class} is a tuple
\[
E(f)=(k,f,P,\Gamma,\Pi,\omega)
\]
containing one class key $k$, one covered public-anchor family $f$, a minimal public-anchor path set $P$, a canonical support-path set $\Gamma$, a maintained packet-key set $\Pi$, and one explanatory sentence $\omega$.
The class does not create a new answer object.
It only gives a stable name to one omitted-support family that already resolves to maintained paths.
\end{definition}

\begin{definition}[Family-coverage rule]
A requestable evidence class satisfies the \emph{family-coverage rule} if every public-anchor family it names already exists in the companion public request contract and every minimal public-anchor path it lists is actually shipped or cited for that family. Those repeated family keys remain imported contract vocabulary rather than new ownership of the family basis inside the evidence-class ledger.
So the class stays attached to the public sentence the reader sees first rather than floating as an ungrounded private-bundle label.
\end{definition}

\begin{definition}[Packet-typing rule]
A disclosure packet satisfies the \emph{packet-typing rule} if every evidence-class key it imports resolves to a maintained requestable evidence class whose canonical support paths are included in or implied by that packet's requestable archive slice.
So the packet keeps ownership of the exact slice, while the evidence class contributes only a stable support-family name.
\end{definition}

\begin{definition}[Contract-coverage rule]
A public request contract satisfies the \emph{contract-coverage rule} if every public-anchor family it names imports at least one requestable evidence class and the union of those classes covers the validator companions and the family-specific support slice that make the public sentence auditable.
So a paper-facing contract may stay short while still saying what kind of omitted support stands behind its source-only promise.
\end{definition}
\begin{definition}[Shared-class reuse rule]
A requestable evidence class satisfies the \emph{shared-class reuse rule} if it may be imported by more than one public-anchor family only when its completion semantics remain family-scoped: the class contributes one stable omitted-support name across those families, but it does not imply a standalone certificate that floats above them.
So stable vocabulary may be shared across families without collapsing completion into one archive-wide verdict.
\end{definition}


\begin{definition}[Shared-evidence-vocabulary rule]
A requestable evidence class satisfies the \emph{shared-evidence-vocabulary rule} if its evidence-class key may recur in later public request contracts, fulfillment certificates, bounded response packets, routes, or series-spine mirrors only as imported class vocabulary: those later objects may point back to one stable class definition, but they do not thereby re-own the class's canonical support-path meaning or its public-anchor-family coverage claim.
So later layers may stay terse by repeating the class key without silently becoming the owner of the omitted-support family it names.
\end{definition}

\begin{definition}[Evidence-class sentence normal form]
A requestable evidence class with class key $k$, imported against one covered public-anchor family $f$, exports an \emph{evidence-class sentence normal form} if it supplies exactly one auditable clause naming (i) the class key $k$, (ii) the covered family key $f$, (iii) the canonical omitted-support meaning carried by the class's support-path / packet slice, and (iv) the posture saying that the class remains imported vocabulary with family-scoped completion rather than a free-standing archive-wide verdict.
\end{definition}

\begin{lemma}[Evidence classes stabilize source-only support promises without widening packet or contract ownership]
If a disclosure packet satisfies the packet-typing rule, every imported requestable evidence class satisfies the shared-class reuse rule when reused across families and the shared-evidence-vocabulary rule when mirrored into later objects, and a public request contract satisfies the family-coverage rule and the contract-coverage rule, then a later terse paper may name one or more requestable evidence classes to characterize its omitted-support promise without changing the ownership of packet paths, review-menu choices, or validator semantics.
\end{lemma}

\begin{proof}
The family-coverage rule binds every evidence class to an already-shipped public-anchor family.
The packet-typing rule prevents evidence classes from inventing a new archive slice; the packet still owns the exact maintained handoff.
The contract-coverage rule then ensures that the paper-facing contract uses those classes only to characterize the omitted-support family that already sits behind its public sentence, while the public-anchor-family keys used to bind those classes still belong to the companion contract layer. The shared-class reuse rule adds the first missing seam: one class may be shared across families, but only as vocabulary imported into those families rather than as a freestanding completion owner. The shared-evidence-vocabulary rule adds the next seam: even when the same class key is mirrored into later certificates, packets, routes, or spine surfaces, the owner of the class's path-family meaning is still the evidence-class ledger rather than those later mirrors.
So evidence classes change only vocabulary: they stabilize support-family names, but they do not replace contract ownership of public-anchor families, packet ownership, menu ownership, fulfillment ownership, or validator ownership.\qedhere
\end{proof}

\section{A minimal public evidence-class card}
\begin{table}[t]
\centering
\small
\begin{tabularx}{0.97\linewidth}{@{}p{0.31\linewidth}X@{}}
\toprule
\textbf{Public row} & \textbf{Pinned value}\\
\midrule
Worked ledger floor & \nolinkurl{example_requestable_evidence_classes.json} with three worked class keys rather than one prose-only omitted-support paragraph \\
Shared class floor & \texttt{bundle\_identity\_and\_validator\_companions} covers both \texttt{public\_status\_sentence\_family} and \texttt{compact\_diff\_summary\_family} and keeps \texttt{completion\_scope = family\_scoped\_import} \\
Public-status family-local class & \texttt{public\_status\_lineage\_and\_closure\_support} with canonical support paths in \nolinkurl{example_successor_lineage_notice.json}, \nolinkurl{example_successor_claim_support_map.json}, \nolinkurl{example_successor_continuity_verdict.json}, \nolinkurl{example_publication_closure_verdict.json}, and \nolinkurl{example_release_closure_ledger.json} \\
Compact-diff family-local class & \texttt{compact\_diff\_compare\_and\_replay\_support} with canonical support paths in \nolinkurl{example_compare_report.json}, \nolinkurl{example_successor_clause_pack.json}, \nolinkurl{example_successor_continuity_verdict.json}, and \nolinkurl{example_verifier_report.json} \\
Vocabulary posture & Every worked class keeps \texttt{evidence\_class\_vocabulary\_mode = imported\_vocabulary\_class\_owned\_meaning} so later mirrors repeat the key without becoming the owner of its support-family meaning \\
Packet / contract floor & The classes are imported by the worked disclosure packets and by the one paper-facing public request contract rather than replacing those narrower owners \\
Escalation pointer & Widen to Synthesis~56 for the paper-facing promise, Synthesis~58 for family completion, Synthesis~59--63 for paper-level status / notice layers, then widen to Synthesis~73 for the maintained request-side terminal stop, drop to Synthesis~17 for one concrete checked-answer/request-tail seam, and widen to Synthesis~74--75 only for abstract paired-tail / cutover ownership rather than omitted-support typing \\
\bottomrule
\end{tabularx}
\caption{A minimal public evidence-class card. This is the smallest public answer later terse papers can quote directly when the live issue is which stable omitted-support family name one maintained handoff imports and what canonical support-family meaning that name carries.}
\label{tab:minimal-public-evidence-class-card}
\end{table}

This card is intentionally non-decorative: it pins the worked ledger, the shared-versus-family-local split, the class-owned vocabulary posture, and the canonical support-path meaning without silently re-owning the paper-facing contract, the family-scoped completion verdict, the paper-level status token, or the request-side terminal judgment.\cite{series:workedexample}

\paragraph{A tiny public evidence-class sentence card.}
\begin{table}[t]
\centering
\small
\begin{tabularx}{0.97\linewidth}{@{}p{0.31\linewidth}X@{}}
\toprule
\textbf{Clause slot} & \textbf{Pinned value}\\
\midrule
Class key floor & \nolinkurl{bundle_identity_and_validator_companions} from \nolinkurl{example_requestable_evidence_classes.json}\\
Covered family fill & One realized family import such as \nolinkurl{public_status_sentence_family}; the same class may be refilled with \nolinkurl{compact_diff_summary_family} without changing the class-owned meaning\\
Meaning floor & Bundle identity plus validator companions, pinned to \nolinkurl{support_manifest.json}, \nolinkurl{example_artifact_inventory.json}, \nolinkurl{example_validation_report.json}, and \nolinkurl{../tools/validate_example.py}\\
Posture floor & \nolinkurl{evidence_class_vocabulary_mode=imported_vocabulary_class_owned_meaning} together with \nolinkurl{completion_scope=family_scoped_import}\\
Escalation pointer & Stop here only for omitted-support typing; widen to Synthesis~56 for the paper-facing promise, to Synthesis~58 for family completion, to Synthesis~59--63 for paper-level status / notice layers, then widen to Synthesis~73 for the maintained request-side terminal stop, drop to Synthesis~17 for one concrete checked-answer/request-tail seam, and widen to Synthesis~74--75 only for abstract paired-tail / cutover ownership\\
\bottomrule
\end{tabularx}
\caption{A tiny public evidence-class sentence card. This is the smallest public sentence-level rendering later terse papers can quote directly when the live issue is one exact omitted-support class clause rather than the paper-facing promise, the family-scoped completion owner, the later status / wrapper owners, or the request-side terminal tail.}
\label{tab:tiny-public-evidence-class-sentence-card}
\end{table}

This sentence card is intentionally non-decorative: it turns the already-owned omitted-support typing into one exact clause for later terse papers, but it does not re-own the paper-facing promise, the family-scoped completion verdict, the later status / notice owners, or the request-side stop / cutover judgment that neighboring papers already own.\cite{series:workedexample}

\paragraph{A tiny evidence-class sentence fill.}
For the worked shared class, fill the four clause slots with class key \nolinkurl{bundle_identity_and_validator_companions}, one realized family import such as \nolinkurl{public_status_sentence_family}, meaning ``bundle identity and validator companions,'' and posture ``imported vocabulary with family-scoped completion.'' The same fill may then be repeated for \nolinkurl{compact_diff_summary_family} without changing the class-owned meaning or letting completion float above the two families. Later terse papers may quote this one-clause omitted-support typing only while those four anchors stay fixed; if the live issue shifts to what the paper promised on request, widen to Synthesis~56, if it shifts to whether the family cut is complete, widen to Synthesis~58, and if it shifts to the exact outward-facing sentence or selector shown to readers, widen to Synthesis~61--63.\cite{series:workedexample}

\paragraph{A compact first-reopen-owner card.}
Once one maintained evidence-class shortcut goes stale, later terse papers still need one small crosswalk: which owner regains authority \emph{first}? The evidence-class note should answer that directly rather than forcing every later paper to translate one stale omitted-support class clause back into the right paper-facing promise owner, the narrower family-completion owner, the later paper-level status / wrapper owners, or the wider request-side terminal cut from prose.

\begin{table}[t]
\centering
\small
\begin{tabularx}{\linewidth}{@{}p{0.31\linewidth}p{0.23\linewidth}X@{}}
\toprule
Failure class & First reopen owner & Why the evidence-class shortcut is no longer honest \\
\midrule
The same source-only paper and same public-anchor family still exist but the local class key, canonical support-family meaning, or shared-versus-family-local vocabulary posture changes & Reopen the requestable-evidence-class owner in Synthesis~57. & The live issue is still the omitted-support family typing itself: which stable class name, class-owned meaning, and vocabulary posture one maintained handoff imports. Once that local class basis changes, the first honest owner is still this note rather than a wider contract or later status / wrapper note. \\
The same class name still appears nearby but the paper-facing source-only promise, covered public-anchor family basis, or validator-companion promise floor changes & Reopen the public-request-contract owner in Synthesis~56. & Evidence classes import their family basis from the paper-facing promise rather than owning it. Once the contract-owned family coverage or validator floor drifts, the first honest owner is the narrower promise layer rather than the class vocabulary that merely types that promise. \\
The class meaning stays fixed but the family-scoped discharge basis or completion verdict for one promised family changes & Reopen the request-fulfillment-certificate owner in Synthesis~58. & The evidence-class note names what kind of omitted support one family needs, but it does not decide whether that family has actually been discharged. Once the completion verdict or its witness basis changes, the first honest owner is the family-completion certificate. \\
The class basis is fixed and the remaining issue is the paper-level status token, its carryforward, or the visible notice sentence / family / selector shown to readers & Widen to Synthesis~59--63 according to whether the live issue is paper-level status, cross-release carryforward, visible notice wording, canonical notice family, or token-keyed family choice. & At that point the omitted-support typing is no longer the live question. The paper should leave this vocabulary note and widen only to the exact downstream status / wrapper owner whose policy now controls. \\
The class basis and visible wrapper basis are fixed and the remaining issue is only the maintained request-side terminal stop, one concrete checked-answer/request-tail seam, or the abstract paired-tail / fail-closed cutover judgment & Widen to Synthesis~73 for the maintained request-side terminal stop, drop to Synthesis~17 for one concrete seam, and widen to Synthesis~74--75 only for the abstract paired-tail / cutover judgment. & Once the omitted-support family names, completion basis, and visible wrapper basis are already fixed, later terse papers should stop at the maintained request-side terminal, one concrete worked seam, or the abstract paired cutover note rather than reopen this class ledger merely to restate an unchanged class meaning. \\
\bottomrule
\end{tabularx}
\caption{A compact first-reopen-owner card. This is the smallest local map from one stale evidence-class shortcut to the first narrower or wider owner that regains authority.}
\label{tab:first-reopen-owner-card-requestable-evidence-class}
\end{table}

\paragraph{One tiny first-reopen-owner drill.}
If a later terse paper still asks \emph{``what kind of omitted support does this same source-only handoff promise?''} but the paper-facing contract and downstream family certificates stay recognizable while the local class key or canonical support-family meaning changes, Table~\ref{tab:first-reopen-owner-card-requestable-evidence-class} sends the paper back to Synthesis~57 itself. If the class meaning still stands but the covered public-anchor family basis or validator-companion promise floor changes, the same table instead reopens Synthesis~56. If the class basis remains fixed while the compact-diff family completion verdict changes, it reopens Synthesis~58. If the class basis is fixed and the real question becomes the paper-level status token or the exact outward-facing notice shown to readers, the same table widens directly to Synthesis~59--63 rather than pretending that the evidence-class vocabulary still owns downstream status compression or wrapper selection.\cite{series:workedexample,series:publicrequeststatusenvelopes,series:publicrequestcarryforward,series:publicrequestnotices,series:publicrequestnoticenormalforms,series:publicrequestnoticeselections,series:requestsideterminalcitationnormalform,series:pairedterminalcitationdiscipline,series:pairedterminalcutoverrules}

\section{Worked example pointer}
The maintained worked bundle now ships \nolinkurl{example_requestable_evidence_classes.json}.\cite{series:workedexample}
It defines three stable evidence classes for the worked note: one shared cross-family class for bundle identity and validator companions, one family-local class for the public-status lineage and closure slice, and one family-local class for the compact-diff compare and replay slice. Each class now also exports an explicit evidence-class-vocabulary posture tag, and the worked route/spine surfaces mirror that ledger through a shared \nolinkurl{source_only_evidence_class_mode_map} rather than silently flattening class meaning into later completion or status objects. The companion public request contract now also exports a public-anchor-family vocabulary posture map, and the worked route/spine surfaces mirror that through \nolinkurl{source_only_public_anchor_family_mode_map} so the family keys named by the evidence classes remain visibly contract-owned rather than silently absorbed into the later class ledger.
The disclosure packets and the paper-facing public request contract import those class keys directly, and the companion request-fulfillment certificates then discharge them family-by-family rather than by inventing a standalone bundle-global completion certificate.\cite{series:requestfulfillmentcertificates}
So the source-only worked note can say not only which maintained packet/menu cut stands behind its public sentence and what kind of omitted support that cut promises, but also how contract-owned family keys and evidence-class vocabulary stay reusable while completion remains local to the family actually requested. Under the paired terminal citation discipline and paired cutover rule, later terse papers should usually stop at Synthesis~73 for the maintained request-side terminal stop, drop to Synthesis~17 when they want one concrete checked-answer/request-tail seam, and widen to Synthesis~74--75 only for the abstract route-side stop/widen judgment; they should reopen this class ledger only when the exact omitted-support family names or shared-versus-family-local class boundaries are themselves disputed.\cite{series:requestsideterminalcitationnormalform,series:pairedterminalcitationdiscipline,series:pairedterminalcutoverrules}

\begin{thebibliography}{99}\small

\bibitem{series:artifactpolicy}
Anonymous.
\newblock Artifact and Disclosure Policy for the One-True-Archive (Source-Only Public Release).
\newblock Series synthesis note (Synthesis~6), 2026. (This archive.)

\bibitem{series:challengeanswerdisclosurepackets}
Anonymous.
\newblock Challenge Answer Disclosure Packets for Successor Maintenance: Public Pointers and Requestable Audit Slices under Space Pressure.
\newblock Series synthesis note (Synthesis~52), 2026. (This archive.)

\bibitem{series:challengeanswerreviewmenus}
Anonymous.
\newblock Challenge Answer Review Menus for Successor Maintenance: Smallest-Sufficient Referee Handoffs from Dockets, Profiles, Packets, and Ladders.
\newblock Series synthesis note (Synthesis~54), 2026. (This archive.)

\bibitem{series:publicrequestcontracts}
Anonymous.
\newblock Public Request Contracts for Source-Only Releases: Exact Paper-Level Handoffs from Boilerplate Policy to Maintained Packet Cuts.
\newblock Series synthesis note (Synthesis~56), 2026. (This archive.)


\bibitem{series:requestfulfillmentcertificates}
Anonymous.
\newblock Request Fulfillment Certificates for Source-Only Releases: Exact Completion Criteria for Public Request Contracts and Typed Omitted-Support Cuts.
\newblock Series synthesis note (Synthesis~58), 2026. (This archive.)

\bibitem{series:publicrequeststatusenvelopes}
Anonymous.
\newblock Public Request Status Envelopes for Source-Only Releases: Paper-Level Tokens Compressing Family Completion without Re-owning the Promise.
\newblock Series synthesis note (Synthesis~59), 2026. (This archive.)

\bibitem{series:publicrequestcarryforward}
Anonymous.
\newblock Public Request Carryforward for Source-Only Successor Releases: When Paper-Level Request Status Persists, Advances, or Reopens across Release Evolution.
\newblock Series synthesis note (Synthesis~60), 2026. (This archive.)

\bibitem{series:publicrequestnotices}
Anonymous.
\newblock Public Request Notices for Source-Only Releases: Outward-Facing Sentences for Within-Release Status and Cross-Release Carryforward.
\newblock Series synthesis note (Synthesis~61), 2026. (This archive.)

\bibitem{series:publicrequestnoticenormalforms}
Anonymous.
\newblock Public Request Notice Normal Forms for Source-Only Releases: Canonical Sentence Families for Within-Release Status and Cross-Release Carryforward Notices.
\newblock Series synthesis note (Synthesis~62), 2026. (This archive.)

\bibitem{series:publicrequestnoticeselections}
Anonymous.
\newblock Public Request Notice Selections for Source-Only Releases: Choosing Canonical Notice Families from Within-Release Status and Cross-Release Carryforward Tokens.
\newblock Series synthesis note (Synthesis~63), 2026. (This archive.)

\bibitem{series:requestsideterminalcitationnormalform}
Anonymous.
\newblock Public Request Response Packet Refresh Response Packet Closure Verdicts for Source-Only Successor Releases: Terminal Closure Stops for the Visible Refresh Branch.
\newblock Series synthesis note (Synthesis~73), 2026. (This archive.)

\bibitem{series:pairedterminalcitationdiscipline}
Anonymous.
\newblock Paired Terminal Citation Discipline for Stable Review Spines: One Answer-Side Stop and One Request-Side Capstone for Terse Papers.
\newblock Series synthesis note (Synthesis~74), 2026. (This archive.)

\bibitem{series:pairedterminalcutoverrules}
Anonymous.
\newblock Paired Terminal Cutover Rules for Stable Review Spines: When to Stop at the Checked Answer, When to Widen to Source-Only, and When to Fail Closed to the Reopening Owner.
\newblock Series synthesis note (Synthesis~75), 2026. (This archive.)

\bibitem{series:workedexample}
Anonymous.
\newblock Worked Example: Receipt Interlock for an Anonymous-DHT Reader-Privacy Claim.
\newblock Series synthesis note (Synthesis~17), 2026. (This archive.)

\end{thebibliography}
\end{document}
